% Test de concurrence de mod_stage : imports simultanés, verrous, volumétrie, dans un Moodle réel.
% Compilation : latexmk -pdf test-de-concurrence.tex
\documentclass[11pt,a4paper]{article}
\usepackage[french]{babel}
\usepackage[T1]{fontenc}
\usepackage{lmodern,microtype,booktabs,tabularx,xcolor,hyperref}
\usepackage[margin=2.2cm]{geometry}
\newcommand{\code}[1]{\texttt{\detokenize{#1}}}
\newcommand{\grav}[1]{\textbf{\textcolor{#1}{\textbullet}}}
\title{Test de concurrence des imports de \code{mod_stage}}
\author{Sébastien Lefebvre}
\date{31 octobre 2026}

\begin{document}
\maketitle

\section{Méthode}

Contrairement au test de charge précédent (\code{test-de-charge.tex}), celui-ci tourne dans un
Moodle complet : Moodle~4.5 (\code{MOODLE_405_STABLE}), PHP~8.3, PostgreSQL~16, version du
module issue de la branche \code{main} (commit \code{dc952d1}). Les verrous d'activité sont ceux
de Moodle (verrous consultatifs PostgreSQL).

\begin{itemize}
  \item \textbf{Données.} Une promotion de 1\,500 étudiants, 60 enseignants, 6 comptes DEVE et
    10 thématiques ; une seconde de 300 étudiants pour les parcours HTTP. Jusqu'à 30\,000 stages
    dans une même activité, soit environ vingt fois une promotion réelle.
  \item \textbf{Processus concurrents.} Chaque procédure tourne dans son propre processus PHP
    (\code{worker.php}). Les processus d'une vague démarrent à la même seconde. Une vague est
    suivie d'un contrôle d'invariants : doublons de stages, stages sans plage, lignes orphelines,
    attributions, mobilités, verrous restants, transactions ouvertes.
  \item \textbf{Parcours HTTP.} Les pages d'import (prévisualisation puis confirmation) sont
    appelées par plusieurs sessions DEVE distinctes, sur le serveur intégré de PHP avec dix
    workers et \code{memory_limit = 256M}.
\end{itemize}

Les procédures couvertes sont les suivantes : import CSV des stages et des référents, import
StageVet, import des mobilités d'études, import du suivi historique, restauration d'un export
global, import depuis un autre cours, service web \code{register_entries}, enregistrement en
masse, suppression en masse, tâches planifiées et ad hoc. Les scripts sont dans
\code{docs/audit/concurrence/}.

\section{Synthèse}

\begin{center}
\begin{tabularx}{\linewidth}{@{}l X l@{}}
\toprule
& \textbf{Constat} & \textbf{Type} \\
\midrule
\grav{red} & Import historique confirmé par deux DEVE : 900 stages en double & Bogue \\
\grav{red} & Restauration d'export global confirmée deux fois : 300 stages en double & Bogue \\
\grav{red} & Deux imports de mobilités simultanés : erreur SQL, import interrompu à mi-fichier & Erreur \\
\grav{orange} & Import depuis un autre cours lancé trois fois : thématiques triplées & Bogue \\
\grav{orange} & Import StageVet et suppression en masse : erreur après validation de l'import & Erreur \\
\grav{orange} & Restauration d'un export de 30\,000 stages : mémoire épuisée & Erreur \\
\grav{orange} & Export Excel : 357~s et 1,5~Go pour 30\,000 stages & Ralentissement \\
\grav{orange} & Verrou d'activité tenu plus de 30~s : les imports suivants échouent & Blocage \\
\bottomrule
\end{tabularx}
\end{center}

\noindent\grav{red}~doublons ou perte de données en usage normal ;
\grav{orange}~erreur ou lenteur en cas de forte charge ou de cas particulier.

\section{Bogues et erreurs}

\subsection{Import du suivi historique : doublons à la confirmation}

Deux DEVE prévisualisent le même classeur (900 stages), puis confirment. Les deux confirmations
réussissent et annoncent chacune « 900 stages importés et validés ». L'activité compte alors
1\,800 stages, dont 900 en double.

\textbf{Cause.} Les doublons probables sont calculés à la prévisualisation, puis enregistrés avec
elle. La confirmation (\code{import_historical.php}, lignes 170 à 223) prend bien le verrou de
l'activité, mais elle ne relit pas les stages déjà enregistrés. La concurrence n'est pas
nécessaire pour déclencher le bogue : il suffit que deux prévisualisations soient confirmées
l'une après l'autre, ou qu'un import CSV ou StageVet ait lieu entre la prévisualisation et la
confirmation.

\textbf{Correction proposée.} Sous le verrou, relire \code{stage_get_existing_theme_pairs()} et
écarter les lignes déjà présentes, comme le fait déjà \code{csv_importer::entries()}.

\subsection{Restauration d'un export global : doublons à la confirmation}

Le même scénario, avec l'export d'une promotion de 300 stages, produit 300 doublons.

\textbf{Cause.} \code{global_export_importer::restore()} écrit dans une transaction, mais elle ne
prend pas le verrou de l'activité et ne relit pas les doublons. Le contrôle des doublons
probables n'a lieu qu'à la prévisualisation (\code{import_global.php}, lignes 58 à 75).

\textbf{Correction proposée.} Prendre le verrou dans \code{restore()}, puis écarter les lignes
qui sont devenues des doublons depuis la prévisualisation.

\subsection{Import des mobilités d'études : erreur d'index unique}

Trois imports simultanés de classeurs qui portent sur les mêmes étudiants donnent le résultat
suivant : un import réussit, les deux autres s'arrêtent sur
\code{dml_write_exception} (index unique \code{stageid, userid} de \code{stage_study_mobility}).
L'un d'eux s'arrête après une quarantaine de lignes déjà écrites.

\textbf{Cause.} \code{study_mobility_importer::import()} n'a ni verrou ni transaction.
\code{stage_set_student_study_mobility()} (\code{locallib.php}, ligne 865) lit la ligne existante,
puis l'insère si elle est absente. Entre ces deux étapes, un autre import peut l'avoir créée.

\textbf{Correction proposée.} Exécuter l'import sous le verrou de l'activité et dans une
transaction. À défaut, rattraper l'erreur d'unicité et mettre à jour la ligne à la place, comme
le fait \code{stage_insert_assignment()}.

\subsection{Import depuis un autre cours : thématiques dupliquées}

Trois imports simultanés des thématiques d'une activité source vers une activité vide créent
30 thématiques au lieu de 10, chacune en trois exemplaires. Les modèles de courriel, protégés par
un index unique, ne sont pas dupliqués.

\textbf{Cause.} \code{stage_import_themes()} (\code{locallib.php}, ligne 5343) écarte les
thématiques de même nom d'après une lecture faite au début de la transaction. Aucun verrou ni
aucun index n'empêche deux imports de faire ce contrôle en même temps. Un double clic reste sans
effet, car la session Moodle sérialise les requêtes d'un même utilisateur. Le cas se produit
avec deux comptes, ou deux navigateurs.

\textbf{Correction proposée.} Prendre le verrou de l'activité cible dans
\code{stage_import_from_stage()}.

\subsection{Import StageVet et suppression en masse : erreur après validation}

Un import StageVet de 300 lignes est lancé en même temps qu'une suppression en masse des stages
qu'il met à jour. L'import se termine sur \code{dml_missing_record_exception}. Pourtant, sa
transaction a été validée : 150 stages ont été créés et 150 mis à jour. La DEVE voit une page
d'erreur, et les invitations aux enseignants référents ne sont pas planifiées.

\textbf{Cause.} Après la libération du verrou, la boucle qui prépare les invitations relit
chaque stage avec \code{MUST_EXIST} (\code{csv_importer.php}, ligne 861). La suppression en masse
ne prend pas le verrou : un stage peut donc disparaître dans l'intervalle.

\textbf{Correction proposée.} Lire avec \code{IGNORE_MISSING} et ignorer le stage absent. Il est
aussi possible de faire prendre le verrou à \code{stage_bulk_delete_entries()}.

\subsection{Restauration d'un gros export : mémoire épuisée}

La prévisualisation de la restauration d'un export de 30\,000 stages (classeur de 6~Mo) s'arrête
sur une erreur fatale : \og Allowed memory size of 268435456 bytes exhausted\fg{}
(\code{global_export_importer.php}, ligne 92). La lecture seule du classeur consomme 336~Mo,
soit environ 11~ko par stage. La limite de 256~Mo est donc atteinte vers 15\,000 à
20\,000 stages. Une promotion de 300 stages se lit en 17~Mo.

\textbf{Correction proposée.} Appeler \code{raise_memory_limit(MEMORY_HUGE)} et
\code{core_php_time_limit::raise()} dans \code{import_global.php}, comme le fait déjà
\code{export.php}.

\section{Ralentissements et blocages}

\subsection{Export Excel}

\begin{center}
\begin{tabular}{@{}r r r r@{}}
\toprule
\textbf{Stages} & \textbf{Durée} & \textbf{Requêtes SQL} & \textbf{Mémoire} \\
\midrule
300 & 3,9~s & -- & -- \\
30\,000 & 357~s & 152 & 1\,515~Mo \\
\bottomrule
\end{tabular}
\end{center}

Les requêtes sont déjà groupées : elles représentent 3~s sur 357. Le reste vient de l'écriture
des cellules, à raison de 11~ms par stage, c'est-à-dire une centaine de cellules.
\code{MoodleExcelWorksheet::write_string()} applique deux styles à chaque cellule
(\code{lib/excellib.class.php}, lignes 205 et 456), et c'est là que le temps passe. Pour une
promotion de 1\,500 stages, l'export devrait prendre une vingtaine de secondes. Derrière un
mandataire inverse dont le délai est de 60~s, un export d'environ 5\,000 stages échouerait.

\textbf{Piste.} Écrire la feuille de détail par \code{PhpSpreadsheet} sans style par cellule,
ou en flux (\texttt{\textbackslash core\textbackslash dataformat}). La restauration doit alors
toujours relire les dates.

\subsection{Verrou d'activité}

Le verrou attend au plus 30~s (\code{stage_get_activity_lock()}, \code{locallib.php},
ligne 1236). Un import de 15\,000 lignes le garde environ 23~s. Avec trois gros imports
simultanés, le troisième échoue au bout de 30,2~s sur \og Another registration \dots{} Please
try again in a few seconds\fg{}. Seize appels simultanés du service web, de 1\,500 stages
chacun, passent tout juste : le dernier se termine en 29,85~s. Pendant ce temps, tout
enregistrement de l'activité est bloqué, y compris les formulaires unitaires, l'enregistrement
par l'étudiant et l'attribution des référents.

\textbf{Pistes.} Ajuster le message, qui annonce \og quelques secondes\fg{}, et allonger l'attente
pour les imports. Pour les très gros fichiers, valider la transaction par paquets.

\section{Ce qui tient}

\begin{center}
\begin{tabularx}{\linewidth}{@{}X r r@{}}
\toprule
\textbf{Procédure (1\,500 lignes, seule)} & \textbf{Durée} & \textbf{Requêtes} \\
\midrule
Import CSV des stages & 2,1~s & 4\,544 \\
Import CSV des référents & 1,1~s & 4\,510 \\
Import StageVet (évaluations comprises) & 3,9~s & 9\,017 \\
Import des mobilités & 2,2~s & 3\,004 \\
Service web, lot de 500 stages & 2,3~s & 1\,555 \\
Import StageVet, 15\,000 lignes & 37~s & 90\,054 \\
\bottomrule
\end{tabularx}
\end{center}

\begin{itemize}
  \item Quatre imports CSV simultanés de deux fichiers à moitié recouvrants créent 1\,500 stages,
    sans aucun doublon.
  \item Trois imports simultanés de référents contradictoires donnent exactement deux référents
    par étudiant, sans violation d'index.
  \item Une vague mêlant huit procédures différentes (CSV, StageVet, référents, deux appels du
    service web, deux enregistrements en masse, mobilités) reste cohérente. Les doubles envois
    du service web et de l'enregistrement en masse sont écartés comme doublons.
  \item Suppression en masse et import StageVet, à 0, 300, 800 et 1\,500~ms d'écart : aucune
    ligne orpheline (plages, détails, réponses, check-lists).
  \item Après chaque vague, aucun verrou consultatif ne reste posé et aucune transaction ne
    reste ouverte.
  \item Avec 30\,000 stages, toutes les pages de l'activité répondent en moins de 0,5~s, sauf le
    tableau de bord (1,7~s), la vue de l'activité (0,8~s) et l'export PDF (1,5~s).
  \item Les tâches planifiées sont sérialisées par le cron de Moodle. Les invitations des
    enseignants référents passent par des tâches ad hoc groupées : 30 tâches pour
    15\,000 stages.
\end{itemize}

\end{document}
